
\paragraph{LLM-based Agentic Systems.}
Recent advances in LLM-based agents have demonstrated that combining planning, tool use, memory, and self-reflection can substantially improve performance on complex, multi-step tasks. The reasoning-acting paradigm, pioneered by ReAct \citep{yao2022react}, interleaves chain-of-thought reasoning with grounded environment actions. Test-time search methods such as Tree-of-Thoughts \citep{yao2023treeofthoughts} and Graph-of-Thoughts \citep{besta2023graphofthoughts} further expand the deliberation space through structured exploration. These systems are strengthened by external tool integration, ranging from self-supervised tool selection in Toolformer \citep{schick2023toolformer} to comprehensive tool ecosystems in ToolLLM \citep{qin2023toolllm}, as well as persistent memory mechanisms that maintain context across long horizons \citep{packer2023memgpt}. A critical aspect of robustness is the ability to reflect and self-correct: Reflexion \citep{shinn2023reflexion} introduces verbal self-reflection to revise future actions, while subsequent work explores richer feedback modalities and meta-cognitive strategies for deciding when to reflect versus act. Despite recent advances, the community still lacks a unified view of how agent components interact and fail. \ModelName{} fills this gap with a debugging layer that traces root causes across the plan–act–observe loop, detects failures from tool-use to memory, and generates corrective feedback that enables agents to improve through their mistakes.
\vspace{-0.1in}
\paragraph{Failure Analyses (Single- and Multi-Agent).}
Systematic failure analyses are emerging to quantify where agents break and how errors propagate. In multi-agent settings, recent works examine collaboration and competition dynamics~\citep{zhu2025multiagentbench}, role specialization, and emergent failure modes such as coordination collapse and dialogue drift \citep{cemri2025why,zhang2025which,zhu2025multiagentbench,hyun2025crewwildfire,tran2025multiagent}. For single-agent systems, analyses have highlighted planning brittleness, grounding and tool-use errors, and hallucination-induced cascades \citep{ji2024testing,sung2025verila,ning2024defining}. Compared to these studies, we focus on actionable debugging: moving from descriptive taxonomies to a two-stage detector that both identifies and helps fix root causes.
\vspace{-0.1in}
\paragraph{Test-Time Scaling for Agents.}
Scaling test-time compute improves reasoning and success rates by allocating more deliberation to challenging instances. In agents, this includes tree/graph search in thought space \citep{yao2023treeofthoughts,besta2023graphofthoughts}, best-of-$N$ sampling with self-consistency, and test-time compute scaling strategies \cite{li2025metal}. Our experiments position \ModelName{} orthogonally to compute scaling: even at fixed compute, targeted debugging recovers substantial performance, and when combined with scaling, it channels extra compute to the right failure points.

\vspace{-0.1in}